\section{Introduction}
\label{introduction}
% \vspace{-0.5em}

Thanks to the powerful modeling capabilities of diffusion models, significant progress has been made in open-world visual content generation, as evidenced by numerous foundational text-to-video models~\citep{wang2023modelscope,chen2024videocrafter2}. These models can generate vivid dynamic content based on arbitrary text prompts. However, while text prompts offer flexibility, they fall short of concretely expressing users' intentions, particularly regarding geometric control. 
%
Although existing trajectory control works primarily rely on training ControlNet-like structures~\citep{videocomposer,controlavideo}, we contend that diffusion model itself contains the potential of substantial control over the generated content without necessitating additional training.
%
In this paper, we aim to investigate the dynamics modeling mechanisms of video diffusion models and explore the possibility of explicitly controlling object trajectories by leveraging their internal properties. While most of the existing efforts are made by modifying text embeddings or adjusting attention mechanisms to enable control or editing~\citep{ren2024move,geyer2023tokenflow}, the influence of initial noises on video motion remains under-explored.

For text-to-video diffusion models, there is considerable diversity in the generated content (\textit{i.e.,} motion trajectories) from the same text prompt, depending on the choice of initial noises. This phenomenon motivates us to raise a question: Is it possible to regulate the motion trajectories with some designs over initial noises?
FreeInit~\citep{wu2023freeinit} observed that low-frequency signals are more resistant to additive noises, which makes the diffusion model biased to inherit layout or shape information from the initial noises. Consequently, by arranging the low-frequency components of noises across frames, we can manipulate the inter-frame content correlation, \textit{i.e.,} the temporal movements of the generated video. However, this constraint is not that reliable because the inter-frame region correlation is not directly aligned with object semantics.
Prior works ~\citep{jain2023peekaboo,ma2023trailblazer} have demonstrated that trajectories can also be influenced by adjusting the attention weights assigned to different objects in some specific areas. Thus, to achieve object-level-based trajectory control, we propose to utilize text-based attention to locate the target objects in cooperation with noise space manipulation.

However, introducing alterations to the noise or attention mechanism carries the risk of causing artifacts in the generated videos. For example, applying a local mask to the self-attention operation can cause partially abnormal values because this diverges from the case encountered by the models during training. Furthermore, these minor anomalies can propagate through subsequent layers and become amplified in the following denoising steps, ultimately resulting in the target region being filled with artifacts. We call such phenomenon as \textit{attention isolation}. Previous work~\citep{jain2023peekaboo} suffers from this problem and is easy to generate artifacts in the areas with masks. In our proposed \textbf{FreeTraj} system, we are fully aware of this issue and mitigate these risks by applying our operations to the noise and attention mechanisms with a tailor-made scheme. Instead of hard attention masks used in Peekaboo~\citep{jain2023peekaboo}, our designed soft attention masks relieve the phenomenon of attention isolation. This approach strikes a balance between staying close to the training distribution and maintaining the ability to control trajectories.

In addition, FreeTraj can be seamlessly integrated into the long video generation framework, enriching the motion trajectories within the generated long videos. Current video generation models are typically trained on a restricted number of frames, leading to limitations in generating high-fidelity long videos during inference. FreeNoise~\citep{freenoise} proposes a tuning-free and time-efficient paradigm for longer video generation based on pre-trained video diffusion models. Although FreeNoise brings satisfactory video quality and visual consistency, it has no guarantee for the various trajectories of generated objects, which are supposed to appear in long videos. With the help of some technical points proposed by FreeNoise, our FreeTraj successfully generates trajectory-controllable long videos. FreeTraj is also valuable in larger video generation. When we directly generate videos with resolutions larger than those in the model training process, we will easily get results with duplicated main objects~\cite {he2023scalecrafter}. However, FreeTraj will constrain the information of the main objects to the target areas. Signals of main objects are suppressed in other areas thus the duplication phenomenon will be reduced.

Our contributions are summarized as follows: \textbf{1)} We investigate the mechanism of how initial noises influence the trajectory of generated objects through several instructive phenomenons. \textbf{2)} We propose \textbf{FreeTraj}, an effective paradigm for tuning-free trajectory control with both noise guidance and attention guidance. \textbf{3)} We extend the control mechanism to achieve longer and larger video generation with a controllable trajectory.
